\documentclass[a4paper,12pt]{ctexart}
\usepackage{xcolor}
\usepackage{graphicx}
\usepackage[top=1.8cm, bottom=1.8cm, left=1.8cm, right=1.8cm]{geometry}
\usepackage{array}
\usepackage{hyperref}
\usepackage{booktabs}
\hypersetup{colorlinks=true,linkcolor=blue!70!black}
\linespread{1.35}

\title{\textcolor{blue!70!black}{\Huge\textbf{计算机科学与技术}}}
\subtitle{\textcolor{gray}{信息技术与互联网方向 · 家长版}}
\author{}
\date{}

\begin{document}
\maketitle
\thispagestyle{empty}

\section{核心数据卡}

\begin{tabular}{ll}
\toprule
专业名称 & 计算机科学与技术 \\
专业代码 & 080901 \\
所属门类 & 工学 \\
\midrule
\multicolumn{2}{l}{\textbf{北京选科要求}} \\
\quad 顶尖院校（清北北航） & 物理必选 \\
\quad 其他院校 & 物理必选（极少数不限） \\
\midrule
\multicolumn{2}{l}{\textbf{北京代表院校}} \\
\quad 第一梯队 & 清华计算机、北大信科 \\
\quad 第二梯队 & 北航、北理工、北邮 \\
\quad 第三梯队 & 北交大、北工大 \\
\quad 第四梯队 & 首师大、北联大、北信科 \\
\midrule
\multicolumn{2}{l}{\textbf{毕业三年月薪参考（北京）}} \\
\quad 互联网大厂（字节/腾讯/阿里） & 25K--50K \\
\quad 中型互联网 & 18K--30K \\
\quad 银行/国企IT部门 & 12K--20K \\
\quad 小型/外包公司 & 8K--15K \\
\bottomrule
\end{tabular}

\section{投入产出比估算}

本科四年学费（公办）约2-3万，生活费和杂费约12-20万。
计算机不必须读研，但大厂核心岗位越来越倾向硕士，若读研总投入约30-50万。
毕业后3年累计收入（中位数）约60-100万。
投资回收期约2-4年，是回收最快的专业之一。

\section{需要留意的风险}

\begin{enumerate}
  \item 行业周期性波动——互联网有大小年，2022-2023年的裁员潮就是例子
  \item AI正在改变入门级编程岗位——低代码和AI辅助编程可能减少初级程序员需求
  \item 35岁门槛是真实存在的焦虑——纯执行岗的职业天花板比想象中低
  \item 自学能力强的人可以绕过学历——这意味着孩子毕业后仍需持续和全国自学者竞争
  \item 红利期在收窄——十年前二本计算机也能进大厂，现在985硕士是简历关门槛
\end{enumerate}

\section{给家长的建议}

\textbf{成绩顶尖（清北北航线）}：计算机是最稳妥的选择之一。鼓励本科期间多做实习，尽早确定细分方向（AI/系统/前端/后端等）。如果孩子对科研有兴趣，清北的直博项目值得考虑。

\textbf{成绩优秀（北邮北交大线）}：依然是性价比很高的选择。建议从大一开始刷LeetCode和做开源项目，不要只满足于课内成绩。考研是多数学生的必经之路。

\textbf{成绩中等（首师大北联大线）}：计算机依然能提供高于平均的起薪，但要管理预期——大概率进中小公司或传统企业的IT部门。建议辅修金融/医学信息/数字媒体等方向增加差异化竞争力。

\textbf{数学或英语明显偏弱}：慎重考虑。计算机专业的核心课程（数据结构、算法、机器学习）对数学要求不低，且大量优质学习资源是英文。如果两科都弱，大学四年会比较痛苦。

\end{document}
